\documentclass[11pt]{article}

\marginparwidth 0.5in 
\oddsidemargin 0.25in 
\evensidemargin 0.25in 
\marginparsep 0.25in
\topmargin 0.25in 
\textwidth=6in
\textheight=8in

\begin{document}
\hfill\vbox{\hbox{Jude Shin}
		\hbox{CSC 4888-S01-2268}	
		\hbox{Homework 1}	
		\hbox{\today}}\par

\bigskip
\centerline{\Large\bf RAW Image Processing}\par
\bigskip

\section{A Note on Python Notebooks}

Each code block in the python notebook generally equates to each step.
I developed this in Google Colab in the browser.
The main function on the bottom links all of the transformations together into one function.
This is done so that I don't accidentally write something in the global scope.
From what I understand, all variables are in one global python file.

If you are running in Google Colab, you have to upload the images every time to the local storage.
I also ignored the requirements.txt and just manually installed the libraries using the pip command.

\section{Explanation}

\subsection{Demosaicing}

Before any of the other per-pixel transformations, I performed a demosaicing algorithm on the RAW image. 
The array of floats represent different r g b values in a Bayer pattern.
At each of the points, the rgb values are averaged with particular filters, taking in the channels in questions around the point we are looking at.
Once all of the individual channels are filtered, then I stack all of them on top of each other to get an 3 channel image.

Once I have that rgb image, then I can perform the other simple per pixel transformations for each pixel in the image.

\subsection{White Balance}

I get the mean of each r g b channels, and then multiply each point by a scalar to get the average of each channel to be a particular float.
In our case, we chose 0.25.
Since this is a per pixel transformation, I took advantage of some numpy syntax instead of nested for loops.

\subsection{(Inverse) Gamma Curve}

For each of the values in a channel, I raise the color value to the inverse of a gamma value.
In our case, we chose 0.25.
Since this is a per pixel transformation, I took advantage of some numpy syntax instead of nested for loops.

\subsection{Clip and Quantize}

I used numpy.clip and numpy.round for clipping and converting to 8-bit integers.
I also scaled the values by 255 to convert from the floating points (which we have been computing in) into the integer values which we can use for colors.

\section{Sources Used}

For the clipping part of this homework, I googled the documentation and found that numpy has its own clipping function.
I also ran into some issues at the very end with rounding values, so I googled the documentation for numpy and found that they also had a rounding function.
All other code bits (like indexing and slicing) was referenced from other labs.

\end{document}
